% 夹注组合用例：夹注格高 ×（句读 / 反白 / 专名号 / 书名号）
%
% 为什么有这个文件：功能各自的用例往往只开自己那一项（#82 的 n-char 用例开着
% \无标点模式），两个功能的交叉处没人测。#216（n-char 夹注里句读偏移）就是
% 这样漏掉的。这里把夹注的三种格高与常和夹注同时出现的标记逐一组合。
%
% 从右到左每列：正文 + 夹注，夹注依次为
%   1. 默认格高（与正文相同）
%   2. n-char-per-col = 30（格高比正文小）
%   3. n-char-per-col = 14（格高比正文大）
%   4. grid-height = 30pt
% 每个夹注里都有句读、\反白、\专名号、\书名号；正文里也有，作对照。
%
% test/geometry_test.py 还会对本文件的 PDF 断言：每个句读标记相对所附字
% 的偏移（以字号计）全文一致，不依赖基线图像。
\documentclass{ltc-guji}
\setmainfont{TW-Kai}
\关闭分页
\开启句读模式
\句读设置{句读颜色=red}

\title{夹注组合}
\chapter{夹注组合}

\begin{document}
\begin{正文}
\专名号{肺}氣虛，則鼻息利。\夹注{靈樞本神作盈。\反白{沈刻本}作懣，今從\专名号{東洋}本、\书名号{千金}}少氣。\par
\专名号{肺}氣虛，則鼻息利。\夹注[n-char-per-col=30]{靈樞本神作盈。\反白{沈刻本}作懣，今從\专名号{東洋}本、\书名号{千金}}少氣。\par
\专名号{肺}氣虛，則鼻息利。\夹注[n-char-per-col=14]{靈樞本神作盈。\反白{沈刻本}作懣，今從\专名号{東洋}本、\书名号{千金}}少氣。\par
\专名号{肺}氣虛，則鼻息利。\夹注[grid-height=30pt]{靈樞本神作盈。\反白{沈刻本}作懣，今從\专名号{東洋}本、\书名号{千金}}少氣。\par
\end{正文}
\end{document}
